\article{局部域上中心单代数的尖点非分歧系列}{Cuspidal Unramified Series for Central Simple Algebras over Local Fields}{Paul Gérardin}{1}{157--169}
\origsection{1}{定理的表述（Statement of the theorem）}
\paragraph*{1.1.}设 $A$ 是非阿基米德局部域（nonarchimedean local field）$F$ 上的中心单代数（central simple algebra）。记 $E$ 为 $A$ 的一个非分歧分裂域（unramified splitting field），$\Gamma$ 为 $E/F$ 的 Galois 群。若 $E^\times$ 的拟特征（quasi-character）在 $\Gamma$ 作用下的全部共轭彼此不同，则称其正则（regular）。

对 $F^\times$ 的任一拟特征 $\chi$，分别以 $\chi_{A/F}$ 和 $\chi_{E/F}$ 表示由
\[
\chi_{A/F}=\chi\circ N_{A/F},\qquad \chi_{E/F}=\chi\circ N_{E/F}
\]
定义的 $A^\times$ 和 $E^\times$ 的拟特征，其中 $N_{A/F}$ 是 $A/F$ 的约化范数（reduced norm），$N_{E/F}$ 是 $E/F$ 的范数。

对 $F$ 的任一非平凡特征 $\psi$，$A$ 的特征 $\psi_{A/F}=\psi\circ T_{A/F}$ 和 $E$ 的特征 $\psi_{E/F}=\psi\circ T_{E/F}$ 均非平凡；这里 $T_{A/F}$ 是 $A/F$ 的约化迹（reduced trace），$T_{E/F}$ 是 $E/F$ 的迹。令 $n$ 为 $A$ 的秩（rank），即 $E/F$ 的次数。\footnote{\textbf{译者注（本次校订）：}这里还需规定 $E/F$ 的次数等于中心单代数 $A$ 的次数 $\deg A=\sqrt{\dim_F A}$，并选取 $E$ 到 $A$ 的嵌入；也就是说，$E$ 是非分歧极大子域。任意非分歧分裂域的次数可以更大，不能自动等于 $\deg A$。后文在分裂情形取 $A=\End_FE$，正是采用这一约定。}
\paragraph*{1.2.}有下述定理。
\begin{originaltheorem}{定理（Theorem）}
沿用上述记号，对 $E^\times$ 的任一正则拟特征 $\theta$，可以确定一个 $A^\times$ 的不可约容许表示（irreducible admissible representation）$\theta_{A/F}$，满足：

（a）与两个正则拟特征 $\dot\theta,\ddot\theta$ 相联系的表示 $\dot\theta_{A/F},\ddot\theta_{A/F}$ 等价，当且仅当 $\dot\theta,\ddot\theta$ 在 $\Gamma$ 下共轭；

（b）$\theta_{A/F}$ 在 $A^\times$ 的中心 $F^\times$ 上的限制由拟特征 $x\mapsto(-1)^{(n-1)\operatorname{val}x}\theta(x)$ 给出；

（c）对 $F^\times$ 的任一拟特征 $\chi$，扭曲表示（twisted representation）$\theta_{A/F}\otimes\chi_{A/F}$ 等价于 $(\theta\chi_{E/F})_{A/F}$；

（d）$\theta_{A/F}$ 的逆反表示（contragredient representation）等价于 $(\theta^{-1})_{A/F}$；

（e）$\theta_{A/F}$ 的 $L$ 函数为 $1$；\footnote{\textbf{译者注（本次校订）：}这一断言需排除 $n=1$ 且 $\theta$ 非分歧的情形：此时 $A=E=F$，每个拟特征都“正则”，而 $L(s,\theta)=(1-\theta(\varpi_F)q^{-s})^{-1}$。当 $n>1$ 时，正则性迫使 $\theta$ 在单位群上非平凡，故其 $L$ 因子为 $1$。1.3中“尖点表示的 $L$ 函数为 $1$”同样需要条件；对除代数，模中心紧性使所有不可约表示均为这里意义下的尖点表示，不能据此一概断言 $L=1$。}

（f）$\theta_{A/F}$ 的 $\varepsilon$ 因子为
\[
\varepsilon(\theta_{A/F},\psi)=(-1)^{(n-1)\operatorname{ord}\psi}\varepsilon(\theta,\psi_{E/F}).
\]
此外，当 $A$ 不是除代数时，表示 $\theta_{A/F}$ 是尖点的（cuspidal）。
\end{originaltheorem}
这里，“尖点”是指表示的矩阵系数（coefficients）在模去中心后具有紧支集。
\paragraph*{1.3.}回顾 $L$ 函数和 $\varepsilon$ 因子的定义 [6]。双线性型 $(a,b)\mapsto T_{A/F}(ab)$ 在 $A$ 上非退化；因此，对 $F$ 的任一非平凡加性特征（additive character）$\psi$，映射
\[
a,b\in A\longmapsto\psi_{A/F}(ab)=\psi(T_{A/F}(ab))
\]
给出 $A$ 与其 Pontryagin 对偶群（Pontryagin dual group）之间的同构。$A$ 上相应的自对偶测度（self-dual measure）记为 $d_{A,\psi}$，Schwartz--Bruhat 空间 $\mathcal S(A)$ 上的 Fourier 变换记为 $f\mapsto f_\psi$：
\[
f_\psi(x)=\int_A f(y)\psi_{A/F}(xy)\,d_{A,\psi}y.
\]
令 $\mathcal A(A)$ 为 $A$ 的乘法群 $A^\times$ 的不可约容许表示的等价类集合。

Godement 与 Jacquet 在 [6] 中证明，积分 $\int_{A^\times}f(x)\pi(x)\,dx$，其中 $f\in\mathcal S(A)$、$dx$ 为 $A$ 上任意 Haar 测度、$\pi$ 为等价类属于 $\mathcal A(A)$ 的 $A^\times$ 表示，可通过下列积分引入复参数并作亚纯延拓：
\[
\int_{A^\times}f(x)c(x)|x|_{A/F}^s\,dx
\]
其中 $c$ 是 $\pi$ 的矩阵系数。该积分在 $\operatorname{Re}(s)$ 足够大时收敛，并可亚纯延拓到整个复平面。他们还证明，存在 Euler 因子（Eulerian factor）$L(\pi)$ 和关于 $s$ 为指数函数的因子 $\varepsilon(\pi,\psi)$，使
\[
\frac{\displaystyle\int_{A^\times}f_\psi(a)\pi(a)^{-1}|a|_{A/F}^{-(n-1)/2}\,d_{A,\psi}a}{L(\pi^\vee\otimes|\ |_{A/F})}
=(-1)^{n-m}\varepsilon(\pi,\psi)
\frac{\displaystyle\int_{A^\times}f(a)\pi(a)|a|_{A/F}^{(n+1)/2}\,d_{A,\psi}a}{L(\pi)}.
\]
其中 $\pi^\vee$ 为 $\pi$ 的逆反表示\footnote{\textbf{译者注（本次校订）：}按本节所用的加性自对偶测度，上一公式右侧的幂次应为 $-(n+1)/2$，缺少了负号；左侧的 $-(n-1)/2$ 保持不变。事实上，约化范数的次数为 $n$，加性空间 $A$ 的维数为 $n^2$，Fourier 变换两端的范数幂次之和应为 $-n$；修正后正好满足这一关系。6.1重复使用该公式时也应作相同修正。}，$m$ 为群 $A^\times$ 在 $F$ 上的秩，$|\ |_{A/F}$ 是将 $F$ 上的正规化绝对值与约化范数 $N_{A/F}$ 复合所得的函数。若 $\pi\in\mathcal A(A)$ 是尖点表示，则 $L(\pi)=1$。
\paragraph*{1.4.}下面证明 $A$ 为分裂代数的情形，即 $A$ 是一个 $n$ 维 $F$ 向量空间的自同态代数。除代数情形的证明与此类似，而且更为简单；一般情形则结合了这两种情况。

表示的构造基本沿用 [4] 的方法；此方法可推广到 $p$ 进域上的任意约化群（reductive group）[3]。对 $\GL(n,F)$，T. Shintani 在若干特殊情形 [10]、R. Howe 在一般情形 [7]，都已构造与 $F$ 的 $n$ 次驯分歧扩张（tamely ramified extensions）相联系的尖点表示；两者均假设 $F$ 的剩余特征为奇数。本文的推进在于：去掉了对剩余特征的限制，给出了典范构造，并计算了 $\varepsilon$ 因子。关于中心单代数乘法群中的 Iwahori 子群的正规化子，R. Howe 在 [8] 中给出了一般构造。进一步厘清这里的表示与 G. Lusztig 利用 $\ell$ 进上同调构造的表示 [9] 之间的关系，也是一个值得研究的问题。
\footnotetext{AMS（MOS）学科分类（1970）：主要分类 22E50、20G25。}

\origsection{2}{Weil 群的非分歧诱导表示（Unramified induced representations of the Weil group）}
域 $E$ 的 Weil 群 $W_E$ 的一维容许表示，与 $E^\times$ 的拟特征一一对应。由正合列（exact sequence）
\[
1\longrightarrow E^\times\longrightarrow W_{E/F}\longrightarrow\Gamma\longrightarrow1,
\]
$E^\times$ 的拟特征 $\theta$ 经诱导给出 $W_{E/F}$ 的一个 $n$ 维表示；而 $W_{E/F}$ 是 $W_F$ 的商，因而得到 $W_F$ 的 $n$ 维表示，记为 $\Ind_E^F\theta$。此外，此表示不可约，当且仅当 $\theta$ 在 $\Gamma$ 下的所有变换彼此不同，即 $\theta$ 正则。一般地，令 $\Gamma_0$ 为 $\theta$ 在 $\Gamma$ 中的稳定子，$E_0$ 为 $E$ 中相应的子域。有 $n_0=[E:E_0]$ 个 $E_0^\times$ 的拟特征 $\theta'$ 满足 $\theta=\theta'_{E/E_0}$，且它们各自关于 $F$ 正则；因此，$\Ind_E^F\theta$ 是 $n_0$ 个表示 $\Ind_{E_0}^F\theta'$ 的直和，每个表示均不可约。这证明了下面命题的（a）。对于（e），使用 $L$ 函数的诱导性（inductivity）。对于（f），$\varepsilon$ 因子只对维数为零的虚表示满足诱导不变性。为确定恒等式
\[
\varepsilon(\Ind_E^F\theta,\psi)=\lambda(E/F,\psi)\varepsilon(\theta,\psi_{E/F})
\]
中的因子 $\lambda(E/F,\psi)$，取 $\theta$ 为平凡特征。于是 $\Ind_E^F1$ 为 $\Gamma$ 的正则表示（regular representation），即 $n$ 个非分歧特征 $x\mapsto\zeta^{\operatorname{val}x}$ 之和，其中 $\zeta$ 为 $n$ 个不同的 $n$ 次单位根。因此
\[
\prod_\zeta\varepsilon(\zeta^{\operatorname{val}},\psi)=\lambda(E/F,\psi)\varepsilon(1,\psi_{E/F}).
\]
但若 $q$ 是 $F$ 的剩余域 $\widetilde F$ 的阶，则
\[
\varepsilon(\zeta^{\operatorname{val}},\psi)=(\zeta q^{1/2})^{\operatorname{ord}\psi},\qquad
\varepsilon(1,\psi_{E/F})=q^{n\operatorname{ord}\psi/2}.
\]
由此得到 $\lambda(E/F,\psi)=(-1)^{(n-1)\operatorname{ord}\psi}$。
\begin{originaltheorem}{命题 1（Proposition 1）}
由 $E^\times$ 的拟特征诱导的 $W_F$ 的表示满足：

（a）由两个拟特征 $\dot\theta,\ddot\theta$ 诱导的表示 $\Ind_E^F\dot\theta$、$\Ind_E^F\ddot\theta$ 等价，当且仅当 $\dot\theta,\ddot\theta$ 在 $\Gamma$ 下共轭；

（b）一维表示 $\det\Ind_E^F\theta$ 由拟特征 $x\in F^\times\mapsto(-1)^{(n-1)\operatorname{val}x}\theta(x)$ 给出；

（c）对 $F^\times$ 的任一拟特征 $\chi$，$\Ind_E^F\theta$ 经 $\chi$ 扭曲的表示等价于 $\Ind_E^F(\theta\chi_{E/F})$；

（d）$\Ind_E^F\theta$ 的逆反表示等价于 $\Ind_E^F(\theta^{-1})$；

（e）$L(\Ind_E^F\theta)=L(\theta)$；

（f）$\varepsilon(\Ind_E^F\theta,\psi)=(-1)^{(n-1)\operatorname{ord}\psi}\varepsilon(\theta,\psi_{E/F})$。
\end{originaltheorem}

\origsection{3}{表示的构造（Construction of the representation）}
\paragraph*{3.1.}对 $E^\times$ 的任一拟特征 $\theta$，有 $n$ 个特征 $\theta/\theta^\gamma$，$\gamma\in\Gamma=\Gal(E/F)$。这些商特征在 $F^\times$ 上均平凡。令 $a(\theta,\gamma)=a(\gamma)$ 为 $\theta/\theta^\gamma$ 的导子（conductor）；它是非负整数，满足
\[
\begin{aligned}
a(\gamma)&=a(\gamma^{-1}),&&\gamma\in\Gamma,\\
a(\gamma\gamma')&\leq\max(a(\gamma),a(\gamma')),&&\gamma,\gamma'\in\Gamma,
\end{aligned}
\]
且当 $a(\gamma)\ne a(\gamma')$ 时后一式取等号；对每个 $\gamma\ne1$ 有 $a(\gamma)\geq1$，当且仅当 $\theta$ 正则。

对每个整数 $r\geq0$，令 $\Gamma_r$ 为满足 $a(\gamma)\leq r$ 的元素 $\gamma$ 构成的 $\Gamma$ 的子群，$E_r$ 为它在 $E$ 中的不动子域。记 $A_r=\End_{E_r}E$，则有包含关系
\[
E\subset A_0\subset A_1\subset A_2\subset\cdots\subset A_r\subset\cdots\subset A.
\]
注意，当 $r\geq s\geq\max_{\gamma\in\Gamma}a(\gamma)$ 时，$A_r=A_s$。

代数 $A_r$ 也可视为由 $\Gamma_r$ 在 $E$ 上的作用构成的斜群代数。于是，当 $r\geq s$ 时，有 $A_s$ 线性投影 $T_{A_r/A_s}:A_r\to A_s$，以及 $A_s$ 子模的直和分解
\[
A_r=A_s\oplus A_r^s=\sum_{\gamma\in\Gamma_r/\Gamma_s}\gamma A_s,\qquad
A_r^s=\sum_{\gamma\in\Gamma_r^s}\gamma A_s=\ker T_{A_r/A_s},
\]
其中 $\Gamma_r^s$ 表示 $\Gamma_s$ 在 $\Gamma_r$ 中的补集。此外，由恒等式
\[
T_{A_r/A_s}(xax^{-1})=x(T_{A_r/A_s}a)x^{-1},\quad x\in A_s^\times,\ a\in A_r,\ r\geq s,
\]
群 $A_s^\times$ 通过共轭作用于 $A_r$，并保持此分解的每一项。

$E$ 的理想 $\mathfrak p_E^m$ 所给出的自然滤过（filtration）在每个代数 $A_r$ 上诱导滤过：格 $A_r(\mathfrak p^m)$ 由满足 $a\mathfrak p_E^m\subset\mathcal O_E$ 的 $a\in A_r$ 构成\footnote{\textbf{译者注（本次校订）：}这里的包含条件应为 $a\mathcal O_E\subseteq\mathfrak p_E^m$。按原式，所得格是预期格的负指数版本 $A_r(\mathfrak p^{-m})$，与紧接着“系数取值于 $\mathfrak p_E^m$”的描述及后面的递减滤过不一致。}，其中 $\mathcal O_E$ 为 $E$ 的整数环；它也等于从 $\Gamma_r$ 到 $\mathfrak p_E^m$ 的函数集合。由于 $E/F$ 非分歧，这些滤过与映射 $T_{A_r/A_s}$ 相容。

$E_r$ 的剩余域记为 $\widetilde E_r$，而
\[
\widetilde A_r=A_r(\mathcal O_E)/A_r(\mathfrak p_E)=\End_{\widetilde E_r}\widetilde E.
\]
有限群 $\widetilde A_s^\times$ 通过共轭作用于每个 $\widetilde E_r$ 代数
\[
A_r(\mathfrak p^m)/A_r(\mathfrak p^{m+1})=A_r(\mathfrak p^m/\mathfrak p^{m+1}),
\]
并在 $r\geq s$ 时保持每个子空间 $\gamma A_s(\mathfrak p^m/\mathfrak p^{m+1})$，$\gamma\in\Gamma_r/\Gamma_s$。
\paragraph*{3.2.}现在假设 $\theta$ 为 $E^\times$ 的正则拟特征，即 $A_0=E$。对每个 $r\geq0$，如下定义 $A_r^\times$ 的子群 $K_r$。令 $r'=\lfloor r/2\rfloor$、$r''=\lceil r/2\rceil$，即 $r'$ 为满足 $2r'\leq r$ 的最大整数，$r''$ 为满足 $r\leq2r''$ 的最小整数。于是 $K_0=F^\times$，$K_1$ 是 $A_1^\times$ 中保持 $E$ 内格 $\mathcal O_E$ 的稳定子群，且对 $r>1$，
\[
K_r=1+A_{r-1}(\mathfrak p^{r''})+A_r^{r-1}(\mathfrak p^{r'}).
\]
对每个 $r\geq1$，集合
\[
K_r^+=1+\mathfrak p_E^r+A_r^0(\mathfrak p^{r''})
\]
是 $K_r$ 的正规子群，被每个 $K_s$，$s\leq r$，正规化。此外，$K_s$ 通过共轭作用于 $K_r$；当 $r>s\geq2$ 时，它将 $K_r$ 送入 $K_r^+$。\footnote{\textbf{译者注（本次校订）：}这里应表述为交换子包含关系 $[K_s,K_r]\subseteq K_r^+$，即 $K_s$ 在商群 $K_r/K_r^+$ 上的共轭作用平凡。由于 $K_s$ 已正规化 $K_r$，其共轭作用是 $K_r$ 的自同构，不能把整个 $K_r$ 映入真子群 $K_r^+$。}商群 $K_r/K_r^+$ 如下：
\[
K_1/K_1^+=\widetilde A_1^\times.
\]
当 $r>1$ 时，$K_r/K_r^+$ 为群 $(1+\mathfrak p_E^{r''})/(1+\mathfrak p_E^r)\simeq\mathfrak p_E^{r''}/\mathfrak p_E^r$，除非 $r$ 是奇数且 $A_r\ne A_{r-1}$，此时它是中心扩张（central extension）
\[
0\longrightarrow\mathfrak p_E^{r''}/\mathfrak p_E^r\longrightarrow K_r/K_r^+\longrightarrow A_r^{r-1}(\mathfrak p^{r'}/\mathfrak p^{r''})\longrightarrow0,
\]
由二余循环（2-cocycle）
\[
a,b\in A_r^{r-1}(\mathfrak p^{r'}/\mathfrak p^{r''})\longmapsto T_{A_r/E}(ab)\in\mathfrak p_E^{r-1}/\mathfrak p_E^r
\]
定义。
\paragraph*{3.3.}对每个 $r\geq1$，选取 $E^\times$ 的特征 $\theta_r$，其导子等于 $\max_{\gamma\in\Gamma_r}a(\gamma)$，使 $\theta/\theta_r$ 被群 $\Gamma_r$ 固定；这意味着存在 $E_r^\times$ 的拟特征 $\chi^{(r)}$，使 $\theta=\theta_r\chi_{E/E_r}^{(r)}$。$r=1$ 时，$\theta_1$ 定义一个 $\widetilde E^\times$ 的特征，关于 $\widetilde E_1$ 正则。$r>1$ 时，$\theta_r$ 在 $1+\mathfrak p_E^{r-1}$ 上的限制被 $\Gamma_{r-1}$ 固定，因此存在 $E_{r-1}$ 的特征 $\psi^{(r)}$，使
\[
\theta_r(1+x)=\psi_{E/E_{r-1}}^{(r)}(x),\qquad x\in\mathfrak p_E^{r-1}/\mathfrak p_E^r.
\]
此外，映射
\[
a,b\in A_r^{r-1}(\mathfrak p^{r'}/\mathfrak p^{r''})\longmapsto
\psi_{E/E_{r-1}}^{(r)}(T_{A_r/E}ab)
\]
是一个自对偶配对（self-duality）；其中 $r',r''$ 如 3.2。由此得到以下引理。
\begin{originaltheorem}{引理 1（Lemma 1）}
对每个 $r>1$，存在唯一的 $K_r$ 的不可约表示等价类 $\kappa'_r$，它在 $K_r^+$ 上平凡，且在 $E^\times\cap K_r$ 上的限制为由 $\theta_r$ 给出的等型表示。以
\[
\kappa_r(x)=\chi_{A_r/E_r}^{(r)}(x)\kappa'_r(x),\qquad x\in K_r,
\]
定义 $K_r$ 的表示 $\kappa_r$。则 $\kappa_r$ 的等价类只依赖于 $\theta$，且被任意 $K_s$，$s<r$，的作用固定。
\end{originaltheorem}
当 $r$ 为偶数或 $A_r=A_{r-1}$ 时，结论可直接得到。其余情形仍采用上述 $r',r''$；引入 $K_r/K_r^+$ 的子群 $H(r)$：它是由同一余循环定义的、$A_r^{r-1}(\mathfrak p^{r'}/\mathfrak p^{r''})$ 被 $\mathfrak p_E^{r-1}/\mathfrak p_E^r$ 的中心扩张。它是 Heisenberg 群（Heisenberg group）。若 $K(r)$ 为它在 $K_r$ 中的原像，则 $1+\mathfrak p_E^{r''}$ 正规化 $K(r)$，且 $K_r=(1+\mathfrak p_E^{r''})K(r)$。引理的第一部分由 Heisenberg 群表示的标准性质得到。为证明第二部分，注意 $\kappa'_r$ 在 $K(r)$ 上的限制可通过以下任意群 $H_s(r)$，$1\leq s<r$，分解。
它是 $A_r^{r-1}(\mathfrak p^{r'}/\mathfrak p^{r''})$ 被 $\mathfrak p_{E_s}^{r-1}/\mathfrak p_{E_s}^r$ 的中心扩张，由 $H(r)$ 的中心上的映射 $T_{E/E_s}$ 得到：
\[
0\longrightarrow\mathfrak p_{E_s}^{r''}/\mathfrak p_{E_s}^r\longrightarrow H_s(r)\longrightarrow A_r^{r-1}(\mathfrak p^{r'}/\mathfrak p^{r''})\longrightarrow0,
\]
相应的二余循环为
\[
a,b\longmapsto T_{A_s/E_s}T_{A_r/A_s}(ab)\in\mathfrak p_{E_s}^{r-1}/\mathfrak p_{E_s}^r\subset\mathfrak p_{E_s}^{r''}/\mathfrak p_{E_s}^r.
\]
现在只需注意，$A_s^\times$ 在 $A_r$ 上的共轭作用固定型 $a\mapsto T_{A_s/E_s}T_{A_r/A_s}a$。
\paragraph*{3.4.}进一步，当 $r>s>1$ 时，$K_s$ 在 $K_r$ 上的作用在 $H_s(r)$ 上给出平凡作用。当 $s=1$ 时，群 $H_1(r)$ 是群 $H_1^\gamma(r)$ 的中心积（central product）；后者为子空间
\[
\gamma A_1(\mathfrak p^{r'}/\mathfrak p^{r''})+\gamma^{-1}A_1(\mathfrak p^{r'}/\mathfrak p^{r''})
\]
的原像。群 $H_1^\gamma(r)$ 是此子空间被 $\mathfrak p_{E_1}^{r''}/\mathfrak p_{E_1}^r$ 的中心扩张。若 $\gamma\Gamma_1=\gamma^{-1}\Gamma_1$，其二余循环为
\[
a,b\in\gamma A_1(\mathfrak p^{r'}/\mathfrak p^{r''})\longmapsto T_{A_1/E_1}(ab).
\]
在一般情形下，该余循环为
\[
\begin{aligned}
a,b\in\gamma A_1(\mathfrak p^{r'}/\mathfrak p^{r''})+\gamma^{-1}A_1(\mathfrak p^{r'}/\mathfrak p^{r''})
\longmapsto T_{A_1/E_1}(a'b''+a''b'),
\end{aligned}
\]
其中 $a=a'+a''$、$b=b'+b''$，$a',b'\in\gamma A_1(\mathfrak p^{r'}/\mathfrak p^{r''})$，$a'',b''\in\gamma^{-1}A_1(\mathfrak p^{r'}/\mathfrak p^{r''})$。$K_1$ 在 $K_r$ 上的作用，给出 $K_1/K_1^+=\widetilde A_1^\times$ 在每个子群 $H_1^\gamma(r)$ 上的作用。当 $\gamma\Gamma_1=\gamma^{-1}\Gamma_1$ 时，上述二余循环为 $\widetilde E_1$ 向量空间 $\gamma A_1(\mathfrak p^{r'}/\mathfrak p^{r''})$ 上关于 $\gamma$ 在 $\widetilde E_1$ 中的不动点域 $E'_1$ 的酉型（unitary form），$\widetilde A_1$ 保持此型。由 [5，推论 4.8.2]，$H_1^\gamma(r)$ 的任一不可约表示 $\kappa'_\gamma$，若其中心作用由 $\psi_{E_1/E_{r-1}}^{(r)}$ 给出，就有到 $\widetilde A_1^\times\ltimes H_1^\gamma(r)$ 的典范延拓（canonical extension）。对满足 $\gamma\Gamma_1\ne\gamma^{-1}\Gamma_1$ 的 $\gamma\in\Gamma_r^{r-1}$，群 $\widetilde A_1^\times$ 保持 $H_1^\gamma(r)$ 的商空间中的每个 $\widetilde A_1$ 子模；因此，此群中在 $\mathfrak p_{E_1}^{r-1}/\mathfrak p_{E_1}^r$ 上由 $\psi_{E_1/E_{r-1}}^{(r)}$ 给出的任一不可约表示 $\kappa'_\gamma$，均可典范地延拓到群 $\widetilde A_1\ltimes H_1^\gamma(r)$\footnote{\textbf{译者注（本次校订）：}此处半直积的第一个因子应为单位群 $\widetilde A_1^\times$。所用作用是单位群的共轭作用；代数 $\widetilde A_1$ 本身不构成这里所需的乘法群。}（[5，命题 1.4]；事实上，它在 $\widetilde A_1^\times$ 上的限制的特征为正）。表示 $\kappa'_\gamma$ 的乘积定义引理中的 $K_r$ 表示 $\kappa'_r$，延拓的乘积定义 $\kappa'_r$ 到群 $K_1K_r$ 的延拓 $\widetilde\kappa'_r$。像引理中那样，以 $\chi^{(r)}$ 扭曲 $\widetilde\kappa'_r$ 所得的此群的表示 $\widetilde\kappa_r$，在等价意义下与 $\psi^{(r)}$ 的选择无关。

当 $r=1$ 时，存在唯一的 $K_1$ 的不可约表示等价类 $\kappa'_1$，它在 $K_1^+$ 上平凡，且它与 $\widetilde A_1^\times$ 的 Steinberg 表示的拉回的张量积，是由子群 $\mathcal O_E^\times K_1^+$ 的一维表示 $tx\mapsto\theta_1(t)$，$t\in\mathcal O_E^\times$、$x\in K_1^+$，诱导的表示。它与 [2，定理 8.3] 中的表示 $R_{E_1}(\theta_1)$ 相差因子 $(-1)^{n_1-1}$，其中 $n_1=[E:E_1]$，且它是尖点表示。将 $\kappa_1$ 定义为 $\kappa'_1$ 经 $\chi^{(1)}$ 的扭曲：
\[
\kappa_1(x)=\chi_{A_1/E_1}^{(1)}(x)\kappa'_1(x),\qquad x\in K_1.
\]
其等价类与 $\chi^{(1)}$ 的选择无关。

最后，$K_0=F^\times$ 的表示 $\kappa_0$ 是 $\theta$ 在 $F^\times$ 上的限制。
\begin{originaltheorem}{命题 2（Proposition 2）}
$E^\times$ 的正则拟特征 $\theta$ 通过公式
\[
\kappa_\theta(x_0x_1\cdots x_r\cdots)=\kappa_0(x_0)\otimes\kappa_1(x_1)\otimes\cdots\otimes\widetilde\kappa_r(x_1)\kappa_r(x_r)\otimes\cdots,
\]
典范地定义 $K_\theta$ 的一个不可约表示等价类。\footnote{\textbf{译者注（本次校订）：}此处省略了 $K_\theta$ 的定义：若 $s=\max_{\gamma\in\Gamma}a(\gamma)$，应取 $K_\theta=K_0K_1\cdots K_{\max(1,s)}$。由3.2的正规化关系，这个乘积是子群；更大指标的 $K_r$ 已包含在其中。公式中的张量积只需取相应的有限多个非平凡因子。对 $\widetilde\kappa_r$，严格说是先延拓到 $K_1K_r$，再在公式中使用其对 $K_1$ 的限制。}这里 $x_r\in K_r$，$\kappa_r$ 如上定义；当 $r>1$ 为奇数且 $A_r\ne A_{r-1}$ 时，$\widetilde\kappa_r$ 为其到 $K_1$ 的典范延拓；其他 $r>1$ 时，$\widetilde\kappa_r$ 为平凡表示。
\end{originaltheorem}

\origsection{4}{交换代数（Commuting algebra）}
\paragraph*{4.1.}令 $\pi_\theta$ 为由 $K_\theta$ 的表示 $\kappa_\theta$ 诱导的 $A^\times$ 的表示，其函数空间由 $A^\times$ 上模中心 $F^\times$ 有紧支集的函数构成。其交换代数可等同于从 $A^\times$ 到 $\kappa_\theta$ 的空间的所有映射组成的集合，这些映射在 $K_\theta$ 的左右作用下按表示 $\kappa_\theta$ 变换。\footnote{\textbf{译者注（本次校订）：}这里的函数应取值于 $\End(V_{\kappa_\theta})$，而非表示空间 $V_{\kappa_\theta}$ 本身；它们还须满足相应的模 $K_\theta$ 紧支集条件。两侧的变换律是 $\varphi(k_1gk_2)=\kappa_\theta(k_1)\varphi(g)\kappa_\theta(k_2)$，这也说明函数值必须是线性算子。}
\paragraph*{4.2.}设 $E^\times$ 的拟特征 $\theta$ 满足：对每个 $\gamma\in\Gamma$、$\gamma\ne1$，特征 $\theta/\theta^\gamma$ 的导子均为 $1$。则表示 $\kappa_\theta$ 是由 $\widetilde A^\times$ 的尖点不可约表示 $\kappa'$ 扭曲而来。若 $T$ 是 $A^\times$ 的极大可对角化子群，且 $A=K_\theta TK_\theta$，则交换代数中的函数 $\varphi$ 由它在 $T$ 上的值确定。对不属于 $K_\theta$ 的 $t\in T$，存在 $A^\times$ 的某个真抛物子群的幺幂根基 $U$，使 $t(U\cap K)t^{-1}$ 包含于 $K_\theta\to\widetilde A^\times$ 的核中。于是对 $u\in U\cap K$，
\[
\varphi(tu)=\kappa_\theta(t)\varphi(u)=\varphi(tut^{-1}t)=\varphi(t).
\]
因此，由 $\kappa'$ 的尖点性，$\varphi(t)=0$。\footnote{\textbf{译者注（本次校订）：}前式中 $\kappa_\theta(t)$ 无定义，因为此时 $t\notin K_\theta$。正确的计算是 $\varphi(tu)=\varphi(t)\kappa_\theta(u)=\kappa_\theta(tut^{-1})\varphi(t)=\varphi(t)$。对 $u$ 求平均，再用有限群表示 $\kappa'$ 的尖点性，即得结论。上面的 Cartan 分解也应写作 $A^\times=K_\theta TK_\theta$。}这表明交换代数只含标量，故 $\pi_\theta$ 不可约；由于其矩阵系数模中心有紧支集，表示 $\pi_\theta$ 也为尖点表示。
\paragraph*{4.3.}现在设 $r=\max_{\gamma\in\Gamma}a(\gamma)>1$，$\varphi$ 为表示 $\pi_\theta$ 的交换代数中的非零函数。取 $a\in A^\times$ 使 $\varphi(a)\ne0$。假定已知对某个 $a_r\in A_r^\times$ 有 $a\in K_\theta a_rK_\theta$；由于 $\varphi$ 由其在左右模 $K_\theta$ 的一组代表元上的值确定，可将 $a$ 换成 $a_r$。由 $\kappa_\theta$ 的定义，它限制到 $K_r$ 的子群 $1+A_r(\mathfrak p^{r''})$ 后，由 $\kappa_r$ 的相应限制给出标量作用（$r''$ 为满足 $2r''\geq r$ 的最小整数），即
\[
1+x\longmapsto\chi_{A_r/E_r}^{(r)}(1+x)\theta_r(1+T_{A_r/E}x).
\]
由于 $\varphi(a)\ne0$，对 $x\in A_r(\mathfrak p^{r''})\cap a^{-1}A_r(\mathfrak p^{r''})a$ 有
\[
\theta_r(1+T_{A_r/E}x)=\theta_r(1+T_{A_r/E}(axa^{-1})).
\]
选取 $E$ 的特征 $\alpha$，使 $\alpha(t)=\theta_r(1+t)$ 对 $t\in\mathfrak p_E^{r''}$ 成立。则由 $x\mapsto\alpha_{A_r/E}(axa^{-1})\alpha_{A_r/E}(x)^{-1}$ 定义的 $A_r$ 的特征，在交集 $a^{-1}A_r(\mathfrak p^{r''})a\cap A_r(\mathfrak p^{r''})$ 上平凡；因此可表示为 $\eta(axa^{-1})^{-1}\xi(x)$，其中 $\xi,\eta$ 是 $A_r$ 上在 $A_r(\mathfrak p^{r''})$ 上平凡的两个特征，使对任意 $x\in A_r$ 有
\[
\alpha_{A_r/E}(x)\xi(x)=\alpha_{A_r/E}(axa^{-1})\eta(axa^{-1}).
\]
\paragraph*{4.4.}固定 $E_r$ 的一个阶为 $-r$ 的特征 $\psi$。特征 $\alpha$ 在 $\mathfrak p_E^{r-1}$ 上的限制在 $\mathfrak p_E^r$ 上平凡，且其在 $\Gamma_r$ 中的稳定子为 $\Gamma_{r-1}$，因此存在 $\widetilde E^\times$ 在 $E^\times$ 中的乘法拉回中的元素 $\tau$，使 $E_r(\tau)=E_{r-1}$，且对 $t\in\mathfrak p_E^{r-1}$，有 $\alpha(t)=\psi_{E/E_r}(\tau t)$。
\begin{originaltheorem}{引理 2（Lemma 2）}
沿用这些记号，则：

（a）$A_r$ 的任一在 $A_r(\mathfrak p^{r''})$ 上等于 $\alpha_{A_r/E_r}$ 的特征，都有一个经 $K_r$ 中元素共轭所得的特征，在 $A_{r-1}(\mathfrak p^{r''})+A_r^{r-1}$ 上等于 $x\mapsto\psi_{A_r/E_r}(\tau x)$；

（b）若 $A_r^\times$ 的某个元素将 $A_r$ 的两个在 $A_{r-1}(\mathfrak p^{r''})+A_r^{r-1}$ 上等于 $x\mapsto\alpha_{A_r/E_r}(\tau x)$ 的特征互相共轭，则此元素属于 $A_{r-1}^\times$。
\end{originaltheorem}
证明中利用配对\footnote{\textbf{译者注（本次校订）：}引理2中的加性特征记号需统一：$\alpha$ 是 $E$ 的特征，所以其到 $A_r$ 的拉回应为 $\alpha\circ T_{A_r/E}$，即4.3的 $\alpha_{A_r/E}$，不能写作 $\alpha_{A_r/E_r}$。引理（b）的共同限制应为 $x\mapsto\psi_{A_r/E_r}(\tau x)$，与（a）及后面的应用一致；证明所用的非退化迹配对也可取 $\psi_{A_r/E_r}(ab)$。} $a,b\in A_r\mapsto\alpha_{A_r/E_r}(ab)$ 将 $A_r$ 与其特征群对应起来；$A_r$ 上的作用为共轭。（a）容易证明。对于（b），须证明：若 $b\in A_r^\times$ 满足 $b(\tau+X)b^{-1}=\tau+X'$，其中 $X,X'$ 在 $A_{r-1}(\mathfrak p^{r'})$ 中，$r'+r''=r$，则 $b\in A_{r-1}^\times$。由于 $X,X'$ 中心化 $\tau$，对等式两边取 $q^{nm}$ 次幂，再令 $m\to\infty$，得到 $b\tau b^{-1}=\tau$。因此 $b$ 中心化 $A_{r-1}$ 的中心 $E_{r-1}$，故本身属于 $A_{r-1}^\times$。

现将引理应用于前面的情形。$A_r$ 的特征 $\alpha_{A_r/E_r}\xi$ 与 $\alpha_{A_r/E_r}\eta$，分别被适当的 $K_r$ 中元素共轭为在 $A_{r-1}(\mathfrak p^{r''})+A_r^{r-1}$ 上等于 $x\mapsto\psi_{A_r/E_r}(\tau x)$ 的特征。由于它们原来由 $a\in A_r^\times$ 共轭，引理 2（b）表明 $a\in K_rA_{r-1}^\times K_r$。
\paragraph*{4.5.}由于 $r$ 足够大时 $A_r=A$，由以上逐步化约可知，$\varphi(a)\ne0$ 蕴涵 $a\in KA_1^\times K$。若 $A_1=E$，则 $a\in K_\theta$，故 $\varphi(a)=\kappa_\theta(a)\varphi(1)$，而 $\varphi(1)$ 是表示 $\kappa_\theta$ 上的交织算子（intertwining operator）。$\kappa_\theta$ 不可约，故此算子为标量；表示不可约，且因其矩阵系数模中心有紧支集而为尖点表示。

接着处理 $A_1\ne E$ 的情形。设 $a\in A_1^\times$ 不在 $K_\theta$ 中，即不在 $K_1F^\times$ 中。则存在 $A_1^\times$ 的某个真抛物子群的幺幂根基 $U$，使 $a(U\cap K_1)a^{-1}$ 属于 $K_1^+$，甚至在任意 $r\geq1$ 的 $\kappa_r$ 的核中，因为行列式在任一幺幂元素上平凡。对每个满足 $a(\gamma)=r$、$r>1$ 为奇数的 $\gamma\in\Gamma$，群 $U\cap K_1$ 作用于 $\widetilde E_1$ 向量空间
\[
\gamma A_1(\mathfrak p^{r'}/\mathfrak p^{r''})+\gamma^{-1}A_1(\mathfrak p^{r'}/\mathfrak p^{r''}).
\]
这一作用的不动向量构成一个子空间，可将其提升为 $H_1(\gamma)$ 的一个交换子群 $\widetilde W_\gamma$（记号如 3.4），且 $K_r$ 内存在由幺幂元素构成的子群 $W_\gamma$，其在 $H(r)$ 上的投影为 $\widetilde W_\gamma$。

当 $r$ 为偶数，或没有满足 $a(\gamma)=r$ 的 $\gamma$ 时，定义 $W_r=1$；$r=1$ 时，令 $W_1=U\cap K_1$；在其余情形，令 $W_r$ 为满足 $a(\gamma)=r$ 的 $\gamma$ 所对应的 $W_\gamma$ 按选定次序、无重复地相乘所得的群。于是，对任意 $r$，$aW_ra^{-1}$ 属于 $K_r^+$，且由幺幂元素构成，因而属于 $\kappa_\theta$ 的核。对 $x=w_1\cdots w_r\cdots$、$w_r\in W_r$，有
\[
\varphi(a)=\varphi(ax)=\varphi(a)\kappa_1(w_1)\otimes\cdots\otimes\widetilde\kappa_r(w_1)\kappa_r(w_r)\otimes\cdots.
\]
由下一个引理，算子 $\int_{W_r}\kappa_r(w_r)\,dw_r$ 被 $\widetilde\kappa_r(w_1)$ 固定。因此，以下平均算子
\[
\int_{W_1\cdots W_r\cdots}\kappa_1(w_1)\otimes\cdots\otimes\widetilde\kappa_r(w_1)\kappa_r(w_r)\otimes\cdots\,dw_1\cdots dw_r\cdots
\]
等于
\[
\int_{W_1}\kappa_1(w_1)\,dw_1\otimes\cdots\otimes\int_{W_r}\kappa_r(w_r)\,dw_r\otimes\cdots,
\]
由 $\kappa_1$ 的尖点性，此式为零。

所用引理如下，记号与 [5] 相同。
\begin{originaltheorem}{引理 3（Lemma 3）}
（a）设 $V$ 为有限域上的有限维向量空间，$\eta$ 为 Heisenberg 群 $H(V)$ 的一个在中心上非平凡的不可约表示。令 $U$ 为 $G(V)$ 的某个真抛物子群的幺幂根基，$W$ 为它在 $V\times V^*$ 中的不动点空间。则 $W$ 可提升为 $H(V)$ 中的一个子群，且算子 $\sum_W\eta(w)$ 在 $U$ 通过 Weil 表示 $W^V$ 的作用下固定。

（b）设 $F$ 为有限域上的 $K/k$ 酉向量空间（unitary vector space），$\eta$ 为对应的 Heisenberg 群 $H(F,i)$ 的一个在中心上非平凡的不可约表示。令 $U$ 为酉群 $U(F,i)$ 的某个真抛物子群的幺幂根基，$W$ 为 $F$ 中相应的全各向同性子空间（totally isotropic subspace）。则 $W$ 可提升为 $H(V)$ 中的一个子群，且算子 $\sum_W\eta(w)$ 在 $U$ 通过 Weil 表示 $W^{F,i}$ 的作用下固定。\footnote{\textbf{译者注（本次校订）：}本项（b）中的 Heisenberg 群应为 $H(F,i)$，而非上一项（a）的 $H(V)$；此处表示 $\eta$ 的定义域正是 $H(F,i)$。}
\end{originaltheorem}
（a）的证明容易；对（b），只需采用将 $\eta$ 实现为 $W$ 上函数的方式，其值落在与 Heisenberg 群 $H(W^\perp/W,i)$ 相联系的表示空间中。

综上，得到以下命题。
\begin{originaltheorem}{命题 3（Proposition 3）}
设 $\theta$ 为 $E^\times$ 的正则拟特征。则由 $\kappa_\theta$ 诱导的表示 $\pi_\theta$ 是 $A^\times$ 的尖点不可约容许表示。
\end{originaltheorem}

\origsection{5}{交织算子（Intertwining operators）}
\paragraph*{5.1.}现在取 $\dot\theta,\ddot\theta$ 为 $E^\times$ 的两个正则拟特征。与 $\dot\theta$、$\ddot\theta$ 对应的对象，分别用一点、两点标记，如 $\dot\pi=\pi_{\dot\theta}$、$\dot K=K_{\dot\theta}$、$\dot\kappa=\kappa_{\dot\theta}$ 等。若一个对象对 $\dot\theta$、$\ddot\theta$ 相同，就省略点号。表示 $\dot\pi,\ddot\pi$ 之间的交织算子，可等同于从 $A^\times$ 到 $\dot\kappa$ 的空间至 $\ddot\kappa$ 的空间的线性映射空间中的映射 $\varphi$\footnote{\textbf{译者注（本次校订）：}为与下面的左右变换公式相容，函数值的方向应为 $\varphi(a):V_{\ddot\kappa}\to V_{\dot\kappa}$。这样 $\dot\kappa(x)$ 才能在左侧复合，$\ddot\kappa(y)$ 才能在右侧复合。}，并满足
\[
\varphi(xay)=\dot\kappa(x)\varphi(a)\ddot\kappa(y),\qquad x\in\dot K,\ y\in\ddot K.
\]
\paragraph*{5.2.}若存在 $\gamma\in\Gamma$ 使 $\ddot\theta=\dot\theta^\gamma$，则对任意 $r$，群 $\dot\Gamma_r$ 与 $\ddot\Gamma_r$ 相同。可以取 $\ddot\chi^{(r)}=\dot\chi^{(r)}$、$\ddot\theta_r=\dot\theta_r^\gamma$。此外，对 $r>1$，表示 $\dot\kappa_r^\gamma$ 及其延拓 $\widetilde{\dot\kappa}_r^\gamma$ 具有相应的 $\ddot\kappa_r$ 的所需性质，因此分别等价于 $\ddot\kappa_r$ 与 $\widetilde{\ddot\kappa}_r$。$K_1$ 的表示 $\dot\kappa'_1$ 与 $\ddot\kappa'_1$ 在 $K_1^+$ 上平凡，且它们与 $\widetilde A_1^\times$ 的 Steinberg 表示的拉回的张量积，等价于由 $\mathcal O_E^\times K_1^+$ 的一维表示
\[
tx\longmapsto\ddot\theta_1(t)=\dot\theta_1^\gamma(t),\qquad t\in\mathcal O_E,\ x\in K_1^+,
\]
诱导的表示，因而它们也等价。最后，在 $F^\times$ 上，\footnote{\textbf{译者注（本次校订）：}这里比较的应是 $\dot\kappa_1^{\prime\,\gamma}$ 与 $\ddot\kappa_1'$，而非未经 $\gamma$ 扭曲的 $\dot\kappa_1'$ 与 $\ddot\kappa_1'$；当 $\gamma$ 不属于 $\Gamma_1$ 时，这个扭曲不能省略。这与本段随后得到的 $\ddot\kappa\simeq\dot\kappa^\gamma$ 一致。上一行 $t$ 的取值也应为 $\mathcal O_E^\times$。}$\ddot\theta=\dot\theta$，故 $\ddot\kappa_0=\dot\kappa_0$。于是，群 $\ddot K=\dot K=K$ 的表示 $\ddot\kappa$ 与 $\dot\kappa^\gamma$ 等价，因此 $\ddot\pi$ 与 $\dot\pi^\gamma$ 也等价。但算子 $\dot\pi(\gamma)$ 将表示 $\dot\pi$ 送至 $\dot\pi^\gamma$，故 $\ddot\pi$ 与 $\dot\pi$ 等价。因此，在 $\Gamma$ 下共轭的两个 $E^\times$ 正则拟特征给出的 $A^\times$ 表示等价。
\paragraph*{5.3.}反过来，假设 $\dot\pi,\ddot\pi$ 之间存在非零交织映射 $\varphi$，取 $a\in A^\times$ 使 $\varphi(a)\ne0$。对非负整数 $r$ 考虑性质
\[
P(r):\quad\text{对所有 }t\geq r,\ \dot\Gamma_t=\ddot\Gamma_t,\quad
\text{且存在 }\gamma\in\Gamma\text{ 使 }a\in\dot K A_r^\times\gamma\ddot K.
\]
注意 $P(0)$ 蕴涵 $\ddot K=\dot K$，且 $\varphi(\gamma)$ 交织 $\ddot\kappa$ 与 $\dot\kappa^\gamma$，从而在 $F^\times$ 上 $\ddot\theta=\dot\theta$。在 $1+\mathfrak p_E$ 上，$\kappa_\theta$ 的限制是由 $\theta$ 在 $1+\mathfrak p_E$ 上的限制定义的标量算子，故 $\ddot\theta$ 与 $\dot\theta^\gamma$ 在 $1+\mathfrak p_E$ 上一致。从而得到 $r>1$ 时 $K_r$ 的表示 $\ddot\kappa_r$ 与 $\dot\kappa_r^\gamma$ 以及它们到 $K_1$ 的延拓的等价。同构 $\varphi(\gamma)$ 交织 $\ddot\kappa$ 与 $\dot\kappa^\gamma$ 在 $K_1$ 上的限制。由于 $\ddot\kappa_r$ 限制到 $K_1$ 的特征不为零，得到 $\ddot\kappa_1$ 与 $\dot\kappa_1^\gamma$ 等价且特征相同。由于在 $1+\mathfrak p_E$ 上 $\ddot\chi^{(1)}=\dot\chi^{(1)}$，对应的 $\widetilde A_1^\times$ 表示等价，故 $\ddot\theta_1$ 与 $\dot\theta_1^\gamma$ 在 $\Gamma_1$ 下共轭：$\ddot\theta_1=\dot\theta_1^{\gamma\gamma_1}$，$\gamma_1\in\Gamma_1$。$\Gamma_1$ 在 $F^\times$ 及 $\dot\theta$ 在 $1+\mathfrak p_E$ 上的限制上平凡作用，因此得到 $\ddot\theta=\dot\theta^{\gamma\gamma_1}$：拟特征 $\dot\theta,\ddot\theta$ 在 $\Gamma$ 下共轭。

当 $r$ 足够大时，$P(r)$ 成立；而且，如果 $s\leq r$ 是大于所有 $\gamma\in\Gamma_r$ 的导子 $\dot a(\gamma),\ddot a(\gamma)$ 的任意整数，则 $P(r)$ 就是 $P(s)$。
\paragraph*{5.4.}
\begin{originaltheorem}{引理 4（Lemma 4）}
沿用上述记号，当 $r\geq1$ 时，有蕴涵 $P(r)\Rightarrow P(r-1)$。
\end{originaltheorem}
可假设 $r=\max_{\gamma\in\Gamma}\dot a(\gamma)\geq s=\max_{\gamma\in\Gamma}\ddot a(\gamma)$。先假设 $r\ne1$。由于 $\varphi(a)$ 由双陪集 $\dot Ka\ddot K$ 确定，依性质 $P(r)$，可取 $a\in A_r^\times\Gamma$。令 $r''$、$s''$ 分别为满足 $2r''\geq r$、$2s''\geq s$ 的最小整数。沿用 3.3 的记号。这些整数 $r'',s''$ 大于 $1$；\footnote{\textbf{译者注（本次校订）：}此处应为“不小于 $1$”：例如 $r=2$ 时，$r''=1$；若 $s=1$ 或 $2$，也有 $s''=1$。下文使用 $1+A(\mathfrak p^{r''})$ 等子群时，所需的正是指数至少为 $1$。}$\varphi(a)$ 非零蕴涵在格
\[
L=A_r(\mathfrak p^{s''})\cap a^{-1}A_r(\mathfrak p^{r''})a
\]
上，两个标量算子 $\ddot\kappa(x)$ 与 $\dot\kappa(axa^{-1})$ 相等：
\[
\ddot\theta_r(1+T_{A_r/E}x)\ddot\chi_{A_r/E_r}^{(r)}(1+x)
=\dot\theta_r(1+T_{A_r/E}axa^{-1})\dot\chi_{A_r/E_r}^{(r)}(1+x).
\]
映射 $x\in L\mapsto\dot\chi_{A_r/E_r}^{(r)}(1+x)\ddot\chi_{A_r/E_r}^{(r)}(1+x)^{-1}$ 是一个特征，在 $A_r(\mathfrak p^s)\cap a^{-1}A_r(\mathfrak p^r)a\supset L\cdot L$ 上平凡。因此，此特征通过 $T_{A_r/E_r}$ 分解为 $E_r$ 的一个在 $\mathfrak p_{E_r}^s$ 上平凡的特征 $\zeta$。令 $\dot\alpha$ 为 $E$ 的一个特征，满足对 $t\in\mathfrak p_E^{r''}$，$\dot\alpha(t)=\zeta_{E/E_r}(t)\dot\theta_r(1+t)$；令 $\ddot\alpha$ 为 $E$ 的一个特征，满足对 $t\in\mathfrak p_E^{s''}$，$\ddot\alpha(t)=\ddot\theta_r(1+t)$。$A_r$ 的两个特征 $\ddot\alpha_{A_r/E}$ 与 $x\mapsto\dot\alpha_{A_r/E}(axa^{-1})$ 在 $L$ 上一致，因此存在 $A_r$ 的两个特征 $\dot\xi,\ddot\xi$，其中 $\dot\xi$ 在 $A_r(\mathfrak p^{r''})$ 上平凡、$\ddot\xi$ 在 $A_r(\mathfrak p^{s''})$ 上平凡，满足对所有 $x\in A_r$ 有
\[
\ddot\alpha_{A_r/E}(x)\ddot\xi(x)=\dot\alpha_{A_r/E}(axa^{-1})\dot\xi(axa^{-1}).
\]
现应用引理 2 的（b），右边可写为 $\psi_{A_r/E_r}((\dot\tau+\dot X)baxa^{-1}b^{-1})$，其中某个 $b\in K_r$、$\dot X\in\dot A_{r-1}(\mathfrak p^{r'})$，而 $\dot\tau$ 在 $\widetilde{\dot E}_{r-1}^\times$ 在 $E$ 中的乘法拉回内，且在 $E_r$ 上为本原的（primitive）。左边为某个 $T\in\mathfrak p_E^{-s}$、$X\in A_r(\mathfrak p^{r-s''})$ 所给出的 $\psi_{A_r/E_r}((T+X)x)$。于是有
\[
a^{-1}b^{-1}(\dot\tau+\dot X)ba=T+X.
\]
左边的 $q^{nm}$ 次幂序列存在极限 $a^{-1}b^{-1}\dot\tau ba$，故必有 $r=s$。再次应用引理 2（b），存在 $c\in K_r$ 使
\[
\psi_{A_r/E_r}((T+X)x)=\psi_{A_r/E_r}((\ddot\tau+\ddot X)cxc^{-1}),
\]
其中 $\ddot X\in\ddot A_{r-1}(\mathfrak p^{r'})$，而 $\ddot\tau$ 在 $\widetilde{\ddot E}_{r-1}^\times$ 在 $E$ 中的乘法拉回内，且在 $E_r$ 上本原。对等式
\[
a^{-1}b^{-1}(\dot\tau+\dot X)ba=c^{-1}(\ddot\tau+\ddot X)c
\]
依次取 $q^{nm}$ 次幂，得到 $a^{-1}b^{-1}\dot\tau ba=c^{-1}\ddot\tau c$。这证明 $\dot\tau,\ddot\tau$ 生成 $E$ 的同一子域，因此 $\dot E_{r-1}=\ddot E_{r-1}$，且由 $bac^{-1}$ 给出的共轭是 $E_{r-1}$ 的自同构。由于它是 $F$ 线性的，故由某个 $\gamma\in\Gamma$ 给出；最后，$bac^{-1}\gamma^{-1}$ 固定 $E_{r-1}/E_r$ 的一个本原元，因而属于 $A_{r-1}$。这就完成了 $r>1$ 时的证明。

现设 $r=1$，已知对任意 $r$ 有 $\dot\Gamma_r=\ddot\Gamma_r$；可假设 $\Gamma_1\ne1$。取 $a_1\in A_1^\times$ 使 $a\in Ka_1\gamma K$。如 4.5，若 $a_1\notin K_1F^\times$，则存在子群 $W\subset K$，使 $a_1\gamma W\gamma^{-1}a_1^{-1}$ 包含于 $\dot\kappa$ 的核，且 $\int_W\ddot\kappa(w)\,dw=0$。但这蕴涵 $\varphi(a_1\gamma)=0$。因此 $a_1\in K$，且 $a\in K\Gamma K$；由于两组 $\Gamma_r$ 已知相同，这正是 $P(0)$。

我们已证明如下结果。
\begin{originaltheorem}{命题 4（Proposition 4）}
$E^\times$ 的两个正则拟特征定义同一个 $A^\times$ 的不可约表示等价类，当且仅当它们在群 $\Gamma$ 下共轭。
\end{originaltheorem}

\origsection{6}{ε 因子的计算（Computation of the ε-factors）}
\paragraph*{6.1.}由 $E^\times$ 的正则拟特征 $\theta$ 所定义的表示 $\pi_\theta$，其 $L$ 函数恒为 $1$ [6，命题 5.11]。将1.3中 $\pi_\theta$ 的 $\varepsilon$ 因子公式限制到 $\kappa_\theta$ 等型子空间，得到恒等式
\[
\begin{aligned}
&\int_{K_\theta}f_\psi(x)\kappa_\theta(x)^{-1}|x|_{A/F}^{-(n-1)/2}\,d_{A,\psi}x\\
&\qquad=\varepsilon(\pi_\theta,\psi)\int_{K_\theta}f(x)\kappa_\theta(x)|x|_{A/F}^{(n+1)/2}\,d_{A,\psi}x.
\end{aligned}
\]
这里用到：$\pi_\theta$ 由 $\kappa_\theta$ 诱导且不可约，因而 $\kappa_\theta$ 在 $\pi_\theta$ 限制到 $K_\theta$ 时以重数一出现。

在这一公式中，取 $f$ 为子群 $1+A(\mathfrak p^a)$ 的示性函数，其中整数 $a$ 大于 $\theta$ 的导子。表示 $\kappa_\theta$ 在此群上平凡，公式给出标量算子
\[
\varepsilon(\pi_\theta,\psi)=\int_{K_\theta}\psi_{A/F}(x)\kappa_\theta(x)^{-1}|x|_{A/F}^{-(n-1)/2}\,d_{A,\psi}x.
\]
此处，当 $|x|_{A/F}$ 较大时，积分按主值（principal value）解释，在 $|x|_{A/F}$ 小时按解析延拓（analytic continuation）解释。
\paragraph*{6.2.}设 $\theta$ 为 $E^\times$ 的导子为 $1$ 的正则特征。则 $K_\theta$ 的表示 $\kappa_\theta$ 来自 $K_\theta/K^+=\widetilde A^\times$ 的表示 $\kappa$，故
\[
\varepsilon(\pi_\theta,\psi)=\sum_{K_\theta/K^+}\left(\int_{K^+}\psi_{A/F}(xu)\,d_{A,\psi}u\right)\kappa_\theta(x)^{-1}|x|_{A/F}^{(n+1)/2}.
\]
测度 $d_{A,\psi}$ 赋予 $A(\mathcal O)$ 的体积为\footnote{\textbf{译者注（本次校订）：}本段开头的 $K_\theta/K^+=\widetilde A^\times$ 应改为 $K_1/K^+=\widetilde A^\times$；$K_\theta$ 还含有中心的统一化元幂，故其商另有赋值方向。这里的求和仍须保留这些中心陪集。} $q^{-n^2\operatorname{ord}\psi/2}$，且 $K^+=1+A(\mathfrak p)$。因此，除非 $x\in\mathfrak p^{-\operatorname{ord}\psi-1}A(\mathcal O)$，内积分为零；满足此条件时，它是 $K^+$ 的测度 $q^{-n^2(1+\operatorname{ord}\psi/2)}$。\footnote{\textbf{译者注（本次校订）：}满足该条件时，内积分还要乘以相位 $\psi_{A/F}(x)$。因为写 $u=1+v$ 后，$\psi_{A/F}(xu)=\psi_{A/F}(x)\psi_{A/F}(xv)$；后一个因子在相应格上的积分才等于该格的体积。紧接着的有限和已重新出现这一相位。该有限和中的 $\theta(\pi^{-1-\operatorname{ord}\psi}a)$ 也应为算子 $\kappa_\theta(\pi^{-1-\operatorname{ord}\psi}a)$，因为一般矩阵 $a$ 不在 $E^\times$ 中；若取标量形式，则对算子和取迹并除以 $\dim\kappa_\theta$。}上述表达式成为
\[
\varepsilon(\pi_\theta,\psi)=q^{n\operatorname{ord}\psi/2}q^{-N}
\sum_{a\in\widetilde A^\times}\psi_{A/F}(\pi^{-1-\operatorname{ord}\psi}a)\theta(\pi^{-1-\operatorname{ord}\psi}a)^{-1},
\]
其中 $N=n(n-1)/2$，$\pi$ 为 $E$ 的素元。$\pi_\theta$ 的尖点性蕴涵此和实际上只涉及 $\widetilde A^\times$ 的半单元素，而在这些元素上，$\kappa$ 的特征已知（3.5）。Springer [11] 已作计算，得到
\[
\begin{aligned}
\varepsilon(\pi_\theta,\psi)
&=q^{n\operatorname{ord}\psi/2}(-1)^{n-1}\sum_{t\in\widetilde E^\times}
\psi_{E/F}(\pi^{-1-\operatorname{ord}\psi}t)\theta(\pi^{-1-\operatorname{ord}\psi}t)^{-1}\\
&=(-1)^{n-1}\int_{E^\times}\psi_{E/F}(t)\theta(t)^{-1}\,d_{E,\psi}t\\
&=(-1)^{n-1}\varepsilon(\theta,\psi_{E/F}).
\end{aligned}
\]
\paragraph*{6.3.}现设 $\theta$ 是 $E^\times$ 的导子为偶数 $2m$ 的正则拟特征。计算 $\varepsilon(\pi_\theta,\psi)$ 的积分可按子群 $K'=1+A(\mathfrak p^m)$ 分解：
\[
\varepsilon(\pi_\theta,\psi)=\sum_{K_\theta/K'}
\left(\int_{K'}\psi_{A/F}(xy)\kappa_\theta(y)^{-1}\,d_{A,\psi}y\right)
\kappa_\theta(x)^{-1}|x|_{A/F}^{(n+1)/2}.
\]
选取 $E$ 中的元素 $\theta_\psi$，使 $\theta(1+t)=\psi_{E/F}(\theta_\psi t)$ 对 $t\in\mathfrak p_E^m$ 成立。由 $\kappa_\theta$ 的定义，它在 $K'$ 上的限制由标量
\[
y=1+u\in1+A(\mathfrak p^m)\longmapsto\theta(1+T_{A/E}u)=\psi_{A/F}(\theta_\psi u)
\]
给出。上面的内积分等于 $\psi(x)$ 乘以特征 $\psi_{A/F}((x-\theta_\psi)u)$ 在 $A(\mathfrak p^m)$ 上关于测度 $d_{A,\psi}$ 的积分。因此，除非 $x-\theta_\psi\in\mathfrak p^{-\operatorname{ord}\psi-m}A(\mathcal O)$，它为零；满足此条件时，它是 $K'$ 的测度 $q^{-n^2(m+\operatorname{ord}\psi)}$。\footnote{\textbf{译者注（本次校订）：}这里的测度因子应为 $q^{-n^2(m+\operatorname{ord}\psi/2)}$，与6.2的自对偶测度归一化及下一行公式一致；完整内积分还含有前面已写出的相位 $\psi_{A/F}(x)$。本段最后一行的 $\varepsilon(\pi_\theta,\psi_{E/F})$ 应为 $\varepsilon(\theta,\psi_{E/F})$：此时积分是在 $E^\times$ 上，表示也是 $E^\times$ 的一维表示 $\theta$。}这一条件等价于 $x\in\theta_\psi K'$，于是
\[
\begin{aligned}
\varepsilon(\pi_\theta,\psi)
&=q^{-n^2(m+\operatorname{ord}\psi/2)}|\theta_\psi|_{E/F}^{(n+1)/2}
\psi_{E/F}(\theta_\psi)\theta(\theta_\psi)^{-1}\\
&=|\theta_\psi|_{E/F}^{1/2}\psi_{E/F}(\theta_\psi)\theta(\theta_\psi)^{-1}\\
&=\int_{E^\times}\psi_{E/F}(t)\theta(t)^{-1}\,d_{E,\psi}t
=\varepsilon(\pi_\theta,\psi_{E/F}).
\end{aligned}
\]
\paragraph*{6.4.}现在取 $\theta$ 为 $E^\times$ 的正则拟特征，其导子 $m$ 为大于 $1$ 的奇数。令整数 $m',m''$ 满足 $m''=m'+1$、$2m'<m<2m''$。选取 $\theta_\psi\in E$ 使 $\theta(1+t)=\psi_{E/F}(\theta_\psi t)$ 对 $t\in\mathfrak p_E^{m''}$ 成立。$\theta$ 在 $1+\mathfrak p_E^{m'}$ 上的限制定义 $\mathfrak p_E^{m'}$ 的一个二次特征（second degree character），满足
\[
\theta(1+t)\theta(1+t')=\theta(1+t+t')\psi_{E/F}(\theta_\psi tt').
\]
映射 $t\in\mathfrak p_E^{m'}\mapsto\psi(\theta_\psi t)\theta(1+t)^{-1}$ 在 $\mathfrak p_E^{m''}$ 上平凡，且 $E^\times$ 的拟特征 $\theta$ 的 $\varepsilon$ 因子为
\[
\begin{aligned}
\varepsilon(\theta,\psi_{E/F})
&=\int_{E^\times}\psi_{E/F}(t)\theta(t)^{-1}\,d_{E,\psi}t\\
&=|\theta_\psi|_{E/F}^{1/2}\psi_{E/F}(\theta_\psi)\theta(\theta_\psi)^{-1}
\widehat{\int}_{\mathfrak p_E^{m'}/\mathfrak p_E^{m''}}
\psi_{E/F}(\theta_\psi t)\theta(1+t)^{-1}.
\end{aligned}
\]
这里，对有限集合 $B$，$\widehat{\int}_B$ 表示 $(\operatorname{Card}B)^{1/2}\sum_B$。\footnote{原译注指出：扫描原文的指数确为 $1/2$，通常的正规化有限和使用指数 $-1/2$，正文保留原式。\textbf{译者注（本次校订）：}这里应取 $(\operatorname{Card}B)^{-1/2}\sum_B$。非退化二次特征的 Gauss 和的绝对值为 $\sqrt{\operatorname{Card}B}$，只有这一归一化才能使后面所写的和等于符号 $(-1)^{n-1}$；用正指数会多出因子 $\operatorname{Card}B$。}

接下来，将计算 $\varepsilon(\pi_\theta,\psi)$ 的积分按子群 $K''=1+A(\mathfrak p^{m''})$ 分解：
\[
\varepsilon(\pi_\theta,\psi)=\sum_{K_\theta/K''}
\left(\int_{K''}\psi_{A/F}(xy)\kappa_\theta(y)^{-1}\,d_{A,\psi}y\right)
\kappa_\theta(x)^{-1}|x|_{A/F}^{(n+1)/2}.
\]
令 $K'=1+A(\mathfrak p^{m'})$。除非 $x\in\theta_\psi K'$，内积分为零，故此公式可写成
\[
\varepsilon(\pi_\theta,\psi)=|\theta_\psi|_{E/F}^{1/2}
\widehat{\int}_{\theta_\psi K'/K''}\psi_{A/F}(x)\kappa_\theta(x)^{-1}.
\]
但 $K'/K''$ 同构于
\[
A(\mathfrak p^{m'}/\mathfrak p^{m''})=\mathfrak p_E^{m'}/\mathfrak p_E^{m''}+A^0(\mathfrak p^{m'}/\mathfrak p^{m''}),
\]
其中 $A^0=\ker T_{A/E}$。于是，上述表达式分解为两个因子。利用 $\varepsilon(\theta,\psi_{E/F})$ 的公式，得到
\[
\varepsilon(\pi_\theta,\psi)=\varepsilon(\theta,\psi_{E/F})
\widehat{\int}_{A^0(\mathfrak p^{m'}/\mathfrak p^{m''})}
\kappa_\theta(1+x)^{-1}\theta(\theta_\psi)\kappa_\theta(\theta_\psi)^{-1}.
\]
令 $r$ 为所有 $\gamma\ne1$ 的 $\theta/\theta^\gamma$ 的导子 $a(\gamma)$ 的最大值。比较 $\theta$ 的导子 $m$ 与这些商特征导子的最大值 $r$，分两种情形讨论。

（a）先假设 $m>r$。如 3.3，有 $\theta=\theta^{(r)}\chi_{E/F}^{(r)}$。$\kappa_\theta$ 在 $K'$ 上的限制为标量，由
\[
1+x\in K'=1+A(\mathfrak p^{m'})\longmapsto
\theta_r(1+T_{A/E}x)\chi_{A/F}^{(r)}(1+x)
\]
给出。对 $x\in A^0(\mathfrak p^{m'})$，它就是 $\chi_{A/F}^{(r)}(1+x)=\chi^{(r)}(1+\sigma_2(x))$，其中 $\sigma_2(x)$ 是将 $\bigwedge^2E$ 视为 $F$ 向量空间时 $x$ 在其上的作用的迹。映射 $x\mapsto\chi^{(r)}(1+\sigma_2(x))$ 是 $\widetilde E$ 上 $(n-1)$ 维向量空间 $A^0(\mathfrak p^{m'}/\mathfrak p^{m''})$ 的非退化二次特征，故其和
\[
\widehat{\int}_{A^0(\mathfrak p^{m'}/\mathfrak p^{m''})}\chi^{(r)}(1+\sigma_2(x))=(-1)^{n-1}.
\]
这给出
\[
\varepsilon(\pi_\theta,\psi)=(-1)^{n-1}\varepsilon(\theta,\psi_{E/F})\theta(\theta_\psi)\kappa_\theta(\theta_\psi)^{-1}.
\]
由表示 $\kappa_\theta$ 的定义，算子 $\theta(\theta_\psi)\kappa_\theta(\theta_\psi)^{-1}$ 只依赖于 $\theta_\psi$ 模 $1+\mathfrak p_E$ 的类。条件 $m>r$ 意味着此类被 $\Gamma$ 固定，故存在 $F^\times$ 中的代表元；在 $F^\times$ 上，$\kappa$ 为 $\theta$，因此 $\kappa_\theta(\theta_\psi)=\theta(\theta_\psi)$。我们证明了公式
\[
\varepsilon(\pi_\theta,\psi)=(-1)^{n-1}\varepsilon(\theta,\psi_{E/F}).
\]

（b）现假设 $m=r$，令 $r'=m'$、$r''=m''$。群 $K'$ 包含于 $K_{r-1}^+K_r$（记号如 3.2）。由 $T_{A/A_{r-1}}$ 定义的分解 $A^0=A_{r-1}^0+A^{r-1}$，对 $x\in A^0(\mathfrak p^{r'})$、$y\in A_{r-1}^0(\mathfrak p^{r'})$、$z\in A^{r-1}(\mathfrak p^{r'})$、$x=y+z$，得到
\[
\kappa_\theta(1+x)=I\otimes\chi_{A_{r-1}/E_{r-1}}^{(r-1)}(1+y)\kappa_r(1+z).
\]
如（a）一样，映射 $y\mapsto\chi_{A_{r-1}/E_{r-1}}^{(r-1)}(1+y)=\chi^{(r-1)}(1+\sigma_2^{r-1}(y))$ 是 $\widetilde E_{r-1}$ 向量空间 $A_{r-1}^0(\mathfrak p^{r'}/\mathfrak p^{r''})$ 上的非退化二次特征，且
\[
\widehat{\int}_{A_{r-1}^0(\mathfrak p^{r'}/\mathfrak p^{r''})}
\chi_{A_{r-1}/E_{r-1}}^{(r-1)}(1+y)^{-1}=(-1)^{n_{r-1}-1},
\]
其中 $n_{r-1}$ 是 $A_{r-1}$ 在 $E_{r-1}$ 上的维数，即 $\Gamma_{r-1}$ 的阶。\footnote{\textbf{译者注（本次校订）：}这里的 $n_{r-1}$ 是中心单代数 $A_{r-1}=\End_{E_{r-1}}E$ 的次数，等于 $[E:E_{r-1}]=|\Gamma_{r-1}|$；其在 $E_{r-1}$ 上的向量空间维数是 $n_{r-1}^2$。}$E$ 的元素 $\theta_\psi$ 模 $1+\mathfrak p_E$ 被 $\Gamma_{r-1}$ 固定，故由表示 $\kappa_\theta$ 的定义（命题 2），算子 $\theta(\theta_\psi)\kappa_\theta(\theta_\psi)^{-1}$ 就是 $I\otimes\widetilde\kappa_r(\theta_\psi)^{-1}$。这意味着 $\varepsilon(\pi_\theta,\psi)$ 与 $\varepsilon(\theta,\psi_{E/F})$ 只相差由下述算子定义的标量：
\[
\widehat{\int}_{A^{r-1}(\mathfrak p^{r'}/\mathfrak p^{r''})}
\kappa_r(1+z)^{-1}\widetilde\kappa_r(\theta_\psi)^{-1}.
\]
在半直积（semidirect product）$H(r)\ltimes\widetilde E^\times$ 中，$\theta_\psi(1+z)$ 的像与 $\theta_\psi$ 的像共轭。这是因为：若 $\gamma\in\Gamma$ 模 $1+\mathfrak p_E$ 固定 $\theta_\psi$，则 $\gamma\in\Gamma_{r-1}$，因此 $z\mapsto\theta_\psi z\theta_\psi^{-1}-z$ 是 $\widetilde E$ 向量空间 $A^{r-1}(\mathfrak p^{r'}/\mathfrak p^{r''})$ 的自同构。为求出上述标量，考察相应算子的迹；此迹就是 $\widetilde\kappa_r(\theta_\psi)^{-1}$ 的迹，等于 $(-1)^{n-n_{r-1}}$（[4，第 16 页定理 1]，或 [5]）。这表明情形（b）中同样有
\[
\varepsilon(\pi_\theta,\psi)=(-1)^{n-1}\varepsilon(\theta,\psi_{E/F}).
\]
\paragraph*{6.5.}前面的计算给出以下结果。
\begin{originaltheorem}{命题 4（Proposition 4，原文编号）}
设 $\theta$ 是 $E^\times$ 的正则特征，$\pi_\theta$ 为第 3 节构造的表示。若 $a(\theta)$ 为 $\theta$ 的导子，则
\[
\varepsilon(\pi_\theta,\psi)=(-1)^{(n-1)a(\theta)}\varepsilon(\theta,\psi_{E/F}).
\]
\end{originaltheorem}
此公式可写为
\[
\varepsilon(\pi_\theta,\psi)=(-1)^{(n-1)\operatorname{ord}\psi}
(-1)^{(n-1)(\operatorname{ord}\psi+a(\theta))}\varepsilon(\theta,\psi_{E/F}).
\]
由 $\theta$ 定义拟特征 $\widetilde\theta=(-1)^{(n-1)\operatorname{val}}\theta$。注意，$E^\times$ 的特征 $(-1)^{(n-1)\operatorname{val}}$，是由 $\Gamma$ 的正则表示的行列式 $\det\Ind_E^F1$ 所定义的 $F^\times$ 的特征到 $E^\times$ 的唯一非分歧延拓。现在可定义表示 $\theta_{A/F}$。

\paragraph*{定义（Definition）}与 $E^\times$ 的正则拟特征 $\theta$ 相联系的表示 $\theta_{A/F}$ 是 $\pi_{\widetilde\theta}$，其中 $\pi_{\widetilde\theta}$ 已在第 3 节定义。

于是，由命题 1、2、3，得到分裂代数 $A=\End_FE$ 的定理。

\origsection{7}{注记（Remark）}
沿用上述记号。对 $E^\times$ 的非正则拟特征 $\theta$，域 $E_0$ 真包含于 $E$，存在 $n_0=[E:E_0]$ 个 $E_0^\times$ 的拟特征 $\theta'$ 满足 $\theta=\theta'_{E/E_0}$，且它们关于 $F$ 正则。由定理，它们定义代数 $A'=\End_FE_0$ 的乘法群的不可约容许表示 $\theta'_{A'/F}$。$n_0$ 个 $A'^\times$ 的直积在 $F$ 上同构于 $A^\times$ 的某个 $F$ 抛物子群 $P$ 的 Levi 子群（Levi subgroup）。由 [1，定理 3]，由所有 $\theta'$ 的表示 $\theta'_{A'/F}$ 的乘积所定义的 $P$ 表示诱导的 $A^\times$ 表示 $\theta_{A/F}$ 不可约。由 [6，定理 3.4]，此表示 $\theta_{A/F}$ 的 $L$ 函数与 $\varepsilon$ 因子分别为各表示 $\theta'_{A'/F}$ 的 $L$ 函数与 $\varepsilon$ 因子的乘积。由诱导性与定理，以 $n'=n/n_0$ 为 $A'$ 的秩，有公式\footnote{\textbf{译者注（本次校订）：}下式中间的 $L$ 因子记号应理解为 $\prod_{\theta'}L(\theta')$，即各 $E_0^\times$ 拟特征的 Tate $L$ 因子的乘积。正确的链式等式是 $L(\theta_{A/F})=\prod_{\theta'}L(\theta'_{A'/F})=\prod_{\theta'}L(\theta')=L(\theta)$；最后一步来自第2节的诱导分解与命题1（e）。原式的 $\theta'_{E_0/F}$ 不符合本文对范数拉回下标的定义。}
\[
L(\theta_{A/F})=\prod L(\theta'_{A'/F})=\prod L(\theta'_{E_0/F})=L(\theta'_{E/E_0})=L(\theta),
\]
以及
\[
\begin{aligned}
\varepsilon(\theta_{A/F},\psi)
&=\prod\varepsilon(\theta'_{A'/F},\psi)
=\prod(-1)^{(n'-1)\operatorname{ord}\psi}\varepsilon(\theta',\psi_{E_0/F})\\
&=(-1)^{n_0(n'-1)\operatorname{ord}\psi}
\varepsilon(\theta'_{E/E_0},\psi_{E/F})\lambda(E/E_0,\psi_{E_0/F})\\
&=(-1)^{n_0(n'-1)\operatorname{ord}\psi+(n_0-1)\operatorname{ord}\psi}
\varepsilon(\theta,\psi_{E/F})\\
&=(-1)^{(n-1)\operatorname{ord}\psi}\varepsilon(\theta,\psi_{E/F}).
\end{aligned}
\]
定理中关于正则拟特征 $\theta$ 的其他函子性（functoriality）性质，也都延伸到这里的情形。

\begin{thebibliography}{99}
\bibitem[1]{c08-1} I. N. Bernstein and A. V. Zelevinskii, \emph{Induced representations of the group GL(n) over a p-adic field}, Functional Anal. Appl. 10 (1976).
\bibitem[2]{c08-2} P. Deligne and G. Lusztig, \emph{Representations of reductive groups over finite fields}, Ann. of Math. (2) 103 (1976), 103--161.
\bibitem[3]{c08-3} P. Gérardin, \emph{On cuspidal representations of p-adic reductive groups}, Bull. Amer. Math. Soc. 81 (1975), 756--758.
\bibitem[4]{c08-4} P. Gérardin, \emph{Construction de séries discrètes p-adiques}, Lecture Notes in Math., vol. 462, Springer, Berlin, 1975.
\bibitem[5]{c08-5} P. Gérardin, \emph{Weil representations associated to finite fields}, J. Algebra 46 (1977), 54--101.
\bibitem[6]{c08-6} R. Godement and H. Jacquet, \emph{Zeta functions of simple algebras}, Lecture Notes in Math., vol. 260, Springer, Berlin, 1970.
\bibitem[7]{c08-7} R. Howe, \emph{Tamely ramified supercuspidal representations of GL\textsubscript{n}}, 预印本。
\bibitem[8]{c08-8} R. Howe, \emph{Tamely ramified base change in local Langlands reciprocity}, 预印本。
\bibitem[9]{c08-9} G. Lusztig, \emph{Some remarks on the supercuspidal representations of p-adic semisimple groups}, 本论文集，第 1 册，第 171--175 页。
\bibitem[10]{c08-10} T. Shintani, \emph{On certain square integrable irreducible unitary representations of some p-adic linear groups}, J. Math. Soc. Japan 20 (1968), 522--505.
\bibitem[11]{c08-11} T. Springer, \emph{The zeta function of a cuspidal representation of a finite group GL\textsubscript{n}(k)}, Lie Groups and Their Representations (Gel'fand, Ed.), Halsted, New York, 1975.
\end{thebibliography}
\par\noindent 巴黎第七大学（Université Paris VII）。
